% !TEX root = ../main.tex
%%%%%%%%%%%%%%%%%%%%%%%%%%%%%%%%%%%%%%%%%%%%%%%%%%%%%%%%%%%%%%%%%%%%%%%
% Discussion
%%%%%%%%%%%%%%%%%%%%%%%%%%%%%%%%%%%%%%%%%%%%%%%%%%%%%%%%%%%%%%%%%%%%%%%

\section{Discussion}
\label{sec:discussion}

\subsection{The second headline: one channel, three falsified hypotheses}

The paper set out to measure a sign frontier and, in the process, found
something more robust than the frontier itself. Three separate stabilisation
hypotheses were posed and each was falsified:

\begin{enumerate}
  \item slower demand expectations would damp the collapse
  (\cref{sec:expectations}) --- they did not;
  \item a balanced-budget demand floor would shrink the collapse region
  (\cref{sec:government}) --- it \emph{enlarged} it, with the fiscal
  instrument saturated at its cap;
  \item that same benefit would cushion firm heterogeneity
  (\cref{sec:heterogeneity}) --- it \emph{lowered} the viability threshold.
\end{enumerate}

All three fail through the same channel and for the same reason. The
instability originates in the wage curve: unemployment moves the wage, the
wage moves profit-maximising employment (increasingly so as $\sigma$ rises),
employment moves unemployment, and each turn of the cycle leaves investment
short of depreciation, so capital erodes. Demand-side instruments act on the
wrong variable. Worse, in case (3) the benefit fails through a mechanism that
is fully traceable and slightly perverse: by \emph{succeeding} at lowering
unemployment, it causes the wage curve to raise the wage, which pushes the
weakest firm below its replacement investment.

The unifying boundary is worth stating plainly, because it is the most
transferable result in the paper: \textbf{demand instruments work where the
binding constraint is demand, and not against wage-driven capital erosion.}

\subsection{What the global analysis does and does not overturn}

The sensitivity analysis is destructive of the project's original headline and
should be read as such: over the defensible parameter space, wage-led is the
exception; $\sigma$ is nearly inert marginally; and viability is governed by a
parameter that the project had classified as a background convention. Any
version of this work that reported $\sigstar = 0.654$ as \emph{the} result
would have been reporting a cell as if it were a space.

What it does not overturn is the conditional finding itself. A conditional
frontier measured at defaults and a marginal effect averaged over fifteen
other parameters are different objects, and \cref{tab:sobol} gives positive
evidence that they should differ here. The honest summary is that the model
has a robust \emph{local} sign structure, an inert \emph{global} one, and an
unresolved reconciliation whose bridging experiment is half-complete
(\cref{sec:bridge}).
